% Source: https://artofproblemsolving.com/community/c3h2209514_can_someone_help_me_with_converting_fractions_into_base_number_decimal_form
% Stable identity: aops-topic-2209514-q01
\textbf{(a)} Since $\dfrac{1}{12}=\dfrac{1}{4}\cdot\dfrac{1}{3}$, first write $\dfrac{1}{3}$ in base $2$. Repeatedly doubling the fractional part: $\dfrac{1}{3}\to\dfrac{2}{3}$ (digit $0$), $\dfrac{2}{3}\to\dfrac{4}{3}$ (digit $1$, keep $\dfrac{1}{3}$), and the pattern repeats, so $\dfrac{1}{3}=0.\overline{01}_2$. (Check: $0.\overline{01}_2=\dfrac{1}{2^2-1}=\dfrac{1}{3}$.) Dividing by $4=2^2$ shifts the point two places:
\[\dfrac{1}{12}=\boxed{0.00\overline{01}_2}=0.000101010\ldots_2.\]

\textbf{(b)} $\left(\dfrac{6}{5}\right)^3=\dfrac{216}{125}=\dfrac{216}{5^3}$. In base $5$, $216=1\cdot125+3\cdot25+3\cdot5+1=1331_5$, and dividing by $5^3$ moves the point three places:
\[\left(\dfrac{6}{5}\right)^3=\boxed{1.331_5}.\]
(Alternatively $\dfrac{6}{5}=1.1_5$ and $(1.1_5)^3$ is computed like $11^3=1331$.)
